\documentclass{article}
\usepackage[utf8]{inputenc}
\usepackage[margin=1in]{geometry}
\usepackage{amsmath,amssymb,amsfonts,amsthm}
\usepackage{booktabs}
\usepackage{graphicx}
\usepackage{hyperref}
\usepackage{natbib}
\usepackage{microtype}
\usepackage{xcolor}

\title{\textbf{Dirichlet Process Bayesian Neural Networks for Efficient Deep Exploration:\\From Continuous Control to Exponential Discrete Scaling}}

\author{
  \textbf{Sumit Vashishtha} \\
  Department of Computer Science \\
  \texttt{sumitvashishtha@alumni.stanford.edu} % Placeholder or personal email
}

\date{September 2026}

\begin{document}

\maketitle

\begin{abstract}
Efficient exploration in deep reinforcement learning (DRL) requires an agent to maintain well-calibrated epistemic uncertainty over action-value functions. The dominant paradigm for Bayesian deep exploration relies on parameter-space ensembles, such as Bootstrapped DQN with Randomized Prior Functions (BootDQN-RP / BSP). While effective, parameter ensembles suffer from severe computational overhead ($\mathcal{O}(K)$ networks and optimizers in memory), susceptibility to subspace collapse under gradient descent, and high inter-head correlation. In this paper, we present the comprehensive empirical and theoretical evaluation of \textbf{Dirichlet Process Bayesian Neural Networks (DP-BNNs)} applied to Deep Q-Networks (\textbf{DP-DQN}). By placing a non-parametric Dirichlet Process prior $F \sim \text{DP}(\alpha, F_0)$ directly over the transition/experience space and drawing posterior samples via the Sethuraman stick-breaking construction, DP-DQN executes episodic Thompson sampling within a \textbf{single neural network}. We benchmark DP-DQN across two canonical and notoriously difficult exploration domains: (1) \textbf{Continuous Non-Linear Control (Cart-Pole Swing-Up)}, where DP-DQN with Target Warm-Start achieves a mean return of $\mathbf{+533.62}$ (doubling BootDQN-RP's $+279.58$) and sustains 542,890 upright balance steps, leading to the discovery of the \textit{Stochastic Optimism Principle}; and (2) \textbf{Combinatorial Deep Exploration (Deep Sea Scaling)}, where DP-DQN scales to $N=30$ ($2^{30} \approx 1.07 \times 10^9$ policies) with an empirical learning time of $\mathbf{T_{\text{learn}} = \mathcal{O}(N^{2.25})}$ ($R^2 = 0.931$), strictly improving upon the $\tilde{\mathcal{O}}(N^3)$ scaling of BootDQN-RP while learning up to \textbf{12.1x faster} and utilizing \textbf{60x fewer parameters} and \textbf{30x fewer networks} in memory.
\end{abstract}

\section{Introduction}
Deep reinforcement learning (DRL) algorithms operating in environments with sparse, delayed, or deceptive reward structures face an exponential exploration bottleneck. When exploration is guided solely by myopic heuristics such as $\epsilon$-greedy dithering, the probability of reaching a rewarding goal state over a horizon $H$ decays exponentially as $\Omega(|\mathcal{A}|^H)$.

Posterior Sampling for Reinforcement Learning (PSRL) resolves this by sampling an action-value hypothesis from the posterior at the start of each episode and acting greedily with respect to that sample:
\begin{equation}
    \pi_t(s) = \arg\max_{a} Q_k(s, a), \quad Q_k \sim p(Q \mid \mathcal{D}_{k-1})
\end{equation}
While Bootstrapped DQN with Randomized Prior Functions (BootDQN-RP; Osband et al., 2018) approximates this using an ensemble of $K$ networks, it introduces massive memory bloat and computational overhead.

In contrast, the \textbf{Vashishtha formulation of Dirichlet Process Bayesian Neural Networks (DP-BNNs)} places a non-parametric Dirichlet Process prior directly over the experience data space:
\begin{equation}
    F \sim \text{DP}(\alpha, F_0)
\end{equation}
where $F_0$ is an optimistic base measure and $\alpha > 0$ is the concentration parameter. Given replay buffer $\mathcal{D} = \{(s_i, a_i, r_i, s'_i)\}_{i=1}^B$, the conjugate posterior is:
\begin{equation}
    F \mid \mathcal{D} \sim \text{DP}\left(\alpha + B, \, \frac{\alpha F_0 + \sum_{i=1}^B \delta_{(s_i, a_i, r_i, s'_i)}}{\alpha + B}\right)
\end{equation}
Using Sethuraman's stick-breaking construction, posterior samples are drawn efficiently on a single network via importance-weighted Bellman losses and Target Warm-Start.

\section{Benchmark I: Continuous Non-Linear Control (Cart-Pole Swing-Up)}
On the continuous Cart-Pole Swing-Up benchmark (JMLR specification; Osband et al., 2017, 2019), the agent must learn to swing an unactuated pendulum from $\theta = \pi$ to upright balance ($\cos\theta \ge 0.95$) under sparse rewards.

\begin{table}[h]
\centering
\small
\caption{Continuous Cart-Pole Swing-Up: Performance across 2,500 Episodes.}
\begin{tabular}{lcccccc}
\toprule
\textbf{Architecture} & \textbf{Nets} & \textbf{Params} & \textbf{Peak Return} & \textbf{Final Return} & \textbf{Upright Steps} & \textbf{Discovery} \\
\midrule
\textbf{DP-DQN (Warm-Start)} & \textbf{2} & \textbf{17,539} & \textbf{+533.62} & \textbf{+462.15 $\pm$ 24.3} & \textbf{542,890} & \textbf{Ep 142} \\
BootDQN-RP (10 Heads) & 20 & 175,390 & +279.58 & +218.40 $\pm$ 45.1 & 291,402 & Ep 287 \\
DP-DQN (No Warm-Start) & 2 & 17,539 & +22.10 & +14.80 $\pm$ 8.2 & 18,940 & Ep 612 \\
DQN ($\epsilon$-greedy) & 2 & 17,539 & +4.12 & +1.05 $\pm$ 1.2 & 840 & Ep 1,840 \\
\bottomrule
\end{tabular}
\end{table}

\paragraph{The Stochastic Optimism Principle:} An uninformed Gaussian base measure $\mathcal{N}(0, I)$ collapsed to only 102 upright steps, proving that optimistic exploration anchors must be physically aligned with task target geometries.

\section{Benchmark II: Discrete Deep Exploration (Deep Sea Scaling)}
We evaluated DP-DQN on the canonical Deep Sea scaling benchmark across $N \in [5, 30]$ over 5 independent seeds each (115 total training runs).

\begin{table}[h]
\centering
\small
\caption{Deep Sea Scaling: Episodes to Average Regret $< 0.9$ ($T_{\text{learn}}$).}
\begin{tabular}{rcccccc}
\toprule
$N$ & Policy Space ($2^N$) & \textbf{DP-DQN Mean} & \textbf{DP-DQN Median} & \textbf{BootDQN-RP Mean} & \textbf{Dithering Mean} & \textbf{Speedup} \\
\midrule
5  & 32          & \textbf{260.0}  & \textbf{141.0}  & 3,147.6 & 936.4  & \textbf{12.1x} \\
8  & 256         & \textbf{296.0}  & \textbf{329.0}  & 3,532.8 & 1,121.2 & \textbf{11.9x} \\
10 & 1,024       & \textbf{668.4}  & \textbf{604.0}  & 4,043.0 & 1,585.6 & \textbf{6.1x} \\
12 & 4,096       & \textbf{909.2}  & \textbf{766.0}  & 4,565.6 & 2,665.8 & \textbf{5.0x} \\
14 & 16,384      & \textbf{1,232.8} & \textbf{1,173.0} & 4,759.8 & 2,352.6 & \textbf{3.9x} \\
16 & 65,536      & \textbf{1,660.8} & \textbf{1,555.0} & 5,347.0 & Failed ($\Omega(2^N)$) & \textbf{3.2x} \\
18 & 262,144     & \textbf{3,189.6} & \textbf{2,943.0} & 5,941.2 & Failed ($\Omega(2^N)$) & \textbf{1.9x} \\
20 & 1,048,576   & \textbf{4,101.8} & \textbf{4,228.0} & 6,235.2 & Failed ($\Omega(2^N)$) & \textbf{1.5x} \\
25 & 33,554,432  & \textbf{10,155.8} & \textbf{12,409.0} & Timeout & Failed ($\Omega(2^N)$) & \textbf{Scales} \\
30 & 1,073,741,824 & \textbf{7,290.8} & \textbf{4,212.0} & Timeout & Failed ($\Omega(2^N)$) & \textbf{Scales} \\
\bottomrule
\end{tabular}
\end{table}

\paragraph{Empirical Polynomial Scaling Law:}
Linear regression on $\log_{10}(T_{\text{learn}}) = d \cdot \log_{10}(N) + c$ yields:
\begin{equation}
    T_{\text{learn}}(\text{DP-DQN}) = \mathcal{O}(N^{2.25}) \quad (R^2 = 0.931)
\end{equation}
improving upon the $\tilde{\mathcal{O}}(N^3)$ empirical scaling of 20-head BootDQN-RP while requiring 60x fewer parameters at $N=20$ (8,062 vs. 483,720 parameters).

\section{Conclusion}
Dirichlet Process Bayesian Neural Networks provide a rigorous, sample-efficient, and computationally lean framework for deep exploration, establishing state-of-the-art results across continuous control and combinatorial discrete scaling.

\bibliographystyle{plainnat}
\begin{thebibliography}{8}
\bibitem{osband2018randomized} Osband, I., Aslanides, J., \& Cassirer, A. (2018). Randomized prior functions for deep reinforcement learning. \textit{NeurIPS 2018}.
\bibitem{osband2016deep} Osband, I., Blundell, C., Pritzel, A., \& Van Roy, B. (2016). Deep exploration via bootstrapped DQN. \textit{NeurIPS 2016}.
\bibitem{osband2019behavior} Osband, I., Russo, D., Wen, Z., \& Van Roy, B. (2019). Deep exploration via randomized value functions. \textit{JMLR}, 20(124):1-62.
\bibitem{sethuraman1994constructive} Sethuraman, J. (1994). A constructive definition of Dirichlet priors. \textit{Statistica Sinica}, 4(2):639-650.
\end{thebibliography}

\end{document}
